% ============================================================================
% Abstract — DUET (NeurIPS 2026)
% Word count: ~211 words (target 180-220)
% Last revised: 2026-05-04 (v2 — rewrite for readability)
% ============================================================================
%
% Design principles for this version:
%   - Open with plain English, not jargon. The 1.5B success-rate numbers
%     come in sentence 1 as a concrete hook.
%   - Each sentence carries ONE idea. No subordinate-clause stacks.
%   - "Exploration is punished." stands alone as a 3-word punch line.
%   - DR3 / BC / SC / Action / State are NOT named in the abstract;
%     they are introduced in the body. Abstract uses plain descriptions.
%   - Em-dash used only once (--- around "experience replay" definition).
%   - Final sentence carries both aggregate (+13pp) and worst-case (+17.5pp).
% ============================================================================

\begin{abstract}
Experience replay from strong teacher agents is an appealing way to overcome the cold-start problem in on-policy RL for LLM agents. Yet in GRPO-style updates, inserting replayed teacher trajectories into the same rollout batch as student trajectories is not a benign form of supervision: it changes both the advantage estimates and the distributional assumptions underlying the policy gradient. We identify two systematic biases induced by this teacher-mixing replay. First, successful teacher rollouts contaminate the group baseline, suppressing the advantages assigned to successful on-policy student rollouts and thereby discouraging exploration. Second, teacher trajectories are off-policy, but exact teacher--student importance correction can be infeasible when teacher likelihoods are unavailable, leaving standard GRPO updates with an uncorrected distribution mismatch. Therefore, we propose \textbf{DUET}, a principled experience-replay framework that first corrects these biases through baseline separation and density-ratio correction, and then extracts complementary teacher signals through two channels: token-level behavior cloning in the action channel and potential-based state-progress shaping in the state channel. Together, these components let DUET exploit teacher trajectories beyond direct replay while adaptively attenuating teacher influence as the student improves, turning teacher data into complementary signals for action imitation and state-guided exploration. Across ALFWorld and WebShop with two model scales, DUET achieves the best performance in all used settings, improving the strongest prior baseline by 13.0 percentage points on average and by 17.5 points in the weakest-agent regime. Code is available at \url{https://anonymous.4open.science/r/DUET-DC26}.
\end{abstract}